\documentclass[11pt]{article}

\usepackage[margin=1in]{geometry}
\usepackage{amsmath,amssymb,amsthm}
\usepackage{algorithm}
\usepackage{algpseudocode}
\usepackage{enumitem}
\usepackage{hyperref}
\usepackage{xcolor}
\usepackage{tikz}
\usepackage{booktabs}

\hypersetup{colorlinks=true, linkcolor=blue, urlcolor=blue, citecolor=blue}

\newtheorem{theorem}{Theorem}[section]
\newtheorem{lemma}[theorem]{Lemma}
\newtheorem{proposition}[theorem]{Proposition}
\newtheorem{corollary}[theorem]{Corollary}
\newtheorem{definition}[theorem]{Definition}
\newtheorem{example}[theorem]{Example}

\title{TOAE209/TOAH209 -- Design and Analysis of Algorithms\\[4pt]
\large Week 9: Dynamic Programming II: Sequence Alignment and Shortest Paths}
\author{Foreign Trade University}
\date{}

\begin{document}
\maketitle
\tableofcontents

\section{Introduction: Two-Dimensional Dynamic Programming}

Last week (Dynamic Programming I) we introduced the dynamic-programming paradigm through
problems whose subproblems were indexed by a \emph{single} parameter -- weighted interval
scheduling (``the best solution using the first $i$ requests'') and the knapsack problem
(``the best value achievable with the first $i$ items and capacity $w$''). This week
(following Kleinberg \& Tardos, Ch.\ 6) we push the same idea further in two directions.

First, we study \textbf{sequence alignment} (a.k.a.\ \emph{edit distance}), where the
subproblems are indexed by \emph{two} parameters -- a prefix of one string and a prefix of
another -- giving a two-dimensional table. This single recurrence underlies the diff
utility, spell checkers, and biological sequence comparison (DNA/protein alignment), and
its close cousin the \textbf{longest common subsequence (LCS)} problem.

Second, we revisit \textbf{shortest paths}, but now in a setting Dijkstra's algorithm
\emph{cannot} handle: graphs with \textbf{negative edge weights}. We will see that
shortest paths have a clean DP formulation -- the \textbf{Bellman--Ford} algorithm -- that
not only tolerates negative edges but can also \emph{detect negative cycles}. We then
generalise to all pairs of vertices with the elegant \textbf{Floyd--Warshall} algorithm,
and finish with the special (and especially fast) case of shortest/longest paths in a
\textbf{directed acyclic graph (DAG)}.

\subsection{The recurring recipe}

Every algorithm this week follows the same four-step DP recipe:
\begin{enumerate}[label=(\arabic*)]
    \item \textbf{Define subproblems.} Identify a small set of subproblems, each indexed by
    one or more parameters, whose optimal solutions can be combined into a solution of the
    whole.
    \item \textbf{Write a recurrence.} Express the optimal value of a subproblem in terms
    of the optimal values of ``smaller'' subproblems, by reasoning about the \emph{last
    decision} an optimal solution must make.
    \item \textbf{Order the computation.} Fill in a table so that every subproblem is
    solved before it is needed (smallest first, or in topological order).
    \item \textbf{Reconstruct.} Trace back through the table (or store ``parent'' pointers)
    to recover an optimal solution itself, not just its value.
\end{enumerate}

The hard step is almost always (2): finding the right way to reason about the last
decision. Once the recurrence is correct, the implementation is mechanical and the running
time is just (number of subproblems) $\times$ (work per subproblem).

\section{Sequence Alignment and Edit Distance}

\subsection{Problem statement}

\begin{definition}[Edit Distance / Sequence Alignment]
Given two strings $X = x_1 x_2 \cdots x_m$ and $Y = y_1 y_2 \cdots y_n$, an
\emph{alignment} is a way of writing the two strings one above the other, inserting
\emph{gap} symbols (``$-$'') into either string, so that the two padded strings have the
same length and no column consists of two gaps. Each column then either
\emph{matches} two equal characters, \emph{mismatches} two unequal characters, or pairs a
character with a gap (an \emph{insertion} or \emph{deletion}). Given a \emph{gap cost}
$\delta \ge 0$ and a \emph{mismatch cost} $\alpha_{pq} \ge 0$ for substituting character
$p$ by $q$ (with $\alpha_{pp} = 0$), the cost of an alignment is the sum of its column
costs, and the goal is to find an alignment of \textbf{minimum total cost}.
\end{definition}

The \textbf{Levenshtein edit distance} is the special case $\delta = 1$ and
$\alpha_{pq} = 1$ for all $p \ne q$: it counts the minimum number of single-character
\emph{insertions}, \emph{deletions}, and \emph{substitutions} needed to turn $X$ into $Y$.
Problem 1 asks you to implement exactly this; Problem 2 generalises to arbitrary $\delta$
and $\alpha$.

\subsection{The structural insight}

Consider an optimal alignment of $X$ and $Y$, and look at its \emph{last} column. There are
exactly three possibilities for what happens to $x_m$ and $y_n$:
\begin{itemize}
    \item $x_m$ and $y_n$ are aligned together (a match if $x_m = y_n$, else a mismatch);
    \item $x_m$ is aligned with a gap ($x_m$ is deleted);
    \item $y_n$ is aligned with a gap ($y_n$ is inserted).
\end{itemize}
Crucially, no other column may involve \emph{both} $x_m$ and $y_n$: alignments preserve
order, so once we have decided what happens to the last characters, the rest of the
alignment is an optimal alignment of the two remaining prefixes. This ``what happens to the
last character'' reasoning is the heart of the recurrence.

\subsection{The recurrence}

Let $\mathrm{OPT}(i,j)$ denote the minimum cost of aligning the prefix $x_1 \cdots x_i$
with the prefix $y_1 \cdots y_j$. Then:
\[
\mathrm{OPT}(i,j) = \min
\begin{cases}
\alpha_{x_i y_j} + \mathrm{OPT}(i-1, j-1) & \text{(align $x_i$ with $y_j$)}\\[2pt]
\delta + \mathrm{OPT}(i-1, j) & \text{(delete $x_i$, gap in $Y$)}\\[2pt]
\delta + \mathrm{OPT}(i, j-1) & \text{(insert $y_j$, gap in $X$)}
\end{cases}
\]
with base cases $\mathrm{OPT}(i, 0) = i\,\delta$ (delete all of $X$'s first $i$ characters)
and $\mathrm{OPT}(0, j) = j\,\delta$ (insert all of $Y$'s first $j$ characters). The answer
is $\mathrm{OPT}(m, n)$.

\begin{theorem}
\label{thm:align}
The recurrence above correctly computes the minimum-cost alignment of $X$ and $Y$.
\end{theorem}

\begin{proof}
We argue that $\mathrm{OPT}(i,j)$ equals the true optimal alignment cost of $x_1\cdots x_i$
and $y_1\cdots y_j$, by induction on $i+j$. The base cases ($i=0$ or $j=0$) are forced:
aligning a non-empty prefix with the empty string requires exactly one gap per character,
at total cost $i\delta$ (resp.\ $j\delta$). For the inductive step with $i,j \ge 1$,
consider any optimal alignment $A$ of the two prefixes and examine its last column. By the
structural insight, that column is one of the three cases above. If it aligns $x_i$ with
$y_j$, then deleting this column leaves an alignment of $x_1\cdots x_{i-1}$ and
$y_1\cdots y_{j-1}$, which must itself be optimal (else we could improve $A$); by induction
its cost is $\mathrm{OPT}(i-1,j-1)$, so $A$ costs $\alpha_{x_i y_j} + \mathrm{OPT}(i-1,j-1)$.
The other two cases are symmetric and give the other two terms. Hence the cost of $A$ is at
least the minimum of the three expressions, so $\mathrm{OPT}(i,j) \le \mathrm{cost}(A)$.
Conversely, each of the three expressions corresponds to a \emph{valid} alignment (take an
optimal alignment of the relevant smaller prefixes and append the appropriate last column),
so $\mathrm{OPT}(i,j)$ is achievable. Therefore $\mathrm{OPT}(i,j)$ equals the optimum.
\end{proof}

\subsection{Running time and the table}

There are $(m+1)(n+1)$ subproblems, each solved in $O(1)$ time from three already-computed
neighbours (the cell directly above, to the left, and diagonally up-left). Filling the
table row by row (or column by column) therefore takes $\mathbf{O(mn)}$ time and
$O(mn)$ space. This is a quadratic-time algorithm for a problem whose naive solution
(enumerating all alignments) is exponential.

\subsection{Reconstructing an optimal alignment}

The value $\mathrm{OPT}(m,n)$ is only the \emph{cost}. To recover the alignment itself
(Problem 3), trace back from cell $(m,n)$ to $(0,0)$: at each cell $(i,j)$, check which of
the three terms attained the minimum, emit the corresponding column (a match/mismatch, a
deletion, or an insertion), and move to the predecessor cell. Reversing the emitted columns
gives the two aligned strings (with gap symbols). The traceback visits at most $m+n$ cells,
so reconstruction is $O(m+n)$ given the filled table.

\begin{algorithm}[h]
\caption{Sequence Alignment (fill table, then trace back)}
\begin{algorithmic}[1]
\State $\mathrm{OPT}[i][0] \gets i\,\delta$ for all $i$; \quad $\mathrm{OPT}[0][j] \gets j\,\delta$ for all $j$
\For{$i = 1$ to $m$}
    \For{$j = 1$ to $n$}
        \State $\mathrm{OPT}[i][j] \gets \min\big(\alpha_{x_i y_j} + \mathrm{OPT}[i-1][j-1],\ \delta + \mathrm{OPT}[i-1][j],\ \delta + \mathrm{OPT}[i][j-1]\big)$
    \EndFor
\EndFor
\State Trace back from $(m,n)$ to $(0,0)$, emitting one column per step
\State \Return $\mathrm{OPT}[m][n]$ and the reconstructed alignment
\end{algorithmic}
\end{algorithm}

\subsection{Hirschberg's space-saving idea}

For very long strings (e.g.\ genomic sequences with $m,n$ in the millions) the $O(mn)$
\emph{space} is the binding constraint, not the time. A beautiful observation rescues us if
we only need the optimal \emph{cost}: each row of the table depends only on the previous
row, so we can keep just two rows (or even one, updated in place) and compute
$\mathrm{OPT}(m,n)$ in $O(\min(m,n))$ space. But then we lose the information needed to
trace back the alignment.

\textbf{Hirschberg's algorithm} recovers the full alignment in $O(mn)$ time and only
$O(m+n)$ space by combining this space-saving trick with divide-and-conquer. The key fact:
an optimal alignment of $X$ and $Y$ must pass through some cell $(m/2, q)$ in the middle
column. The optimal $q$ can be found by computing, in linear space, the cost of aligning the
first half of $X$ with each prefix of $Y$ (a forward pass) and the cost of aligning the
second half of $X$ with each suffix of $Y$ (a backward pass), then choosing the $q$
minimising their sum. The problem then splits into two independent halves
($x_1\cdots x_{m/2}$ vs.\ $y_1\cdots y_q$, and the rest), each solved recursively. The
recurrence for the total work is $T(m,n) = O(mn) + T(m/2, q) + T(m/2, n-q)$, which solves to
$O(mn)$ -- the same asymptotic time, but now in linear space. We mention this only as a
pointer; the details are in Kleinberg \& Tardos, Section 6.7.

\section{Longest Common Subsequence}

\subsection{Problem statement}

\begin{definition}[Longest Common Subsequence]
A \emph{subsequence} of a string is obtained by deleting zero or more characters without
reordering the rest. Given $X = x_1\cdots x_m$ and $Y = y_1\cdots y_n$, the \emph{longest
common subsequence (LCS)} problem asks for a longest string that is a subsequence of both.
\end{definition}

For example, the LCS of \texttt{AGGTAB} and \texttt{GXTXAYB} is \texttt{GTAB} (length 4).

\subsection{The recurrence}

Let $L(i,j)$ be the length of the LCS of $x_1\cdots x_i$ and $y_1\cdots y_j$. Reasoning
again about the last characters:
\[
L(i,j) = \begin{cases}
0 & \text{if } i = 0 \text{ or } j = 0,\\
1 + L(i-1, j-1) & \text{if } x_i = y_j,\\
\max\big(L(i-1, j),\, L(i, j-1)\big) & \text{if } x_i \ne y_j.
\end{cases}
\]
If $x_i = y_j$ we may safely match them (one can prove some optimal LCS does so); otherwise
at least one of $x_i, y_j$ is unused, giving the two $\max$ terms. The table is $O(mn)$ to
fill, and an actual LCS string is recovered by the same kind of traceback as for alignment
(Problem 4). LCS is itself a special case of the alignment framework: it is exactly the
\emph{maximum-score} alignment with match reward $+1$, mismatch $0$, and gap $0$.

\subsection{The LCS--edit-distance relationship}

There is a clean identity when the only edit operations are insertion and deletion (no
substitution -- equivalently, substitution costs $2$, more than a delete-plus-insert):
\begin{theorem}
\label{thm:lcs-edit}
Let $d(X,Y)$ be the minimum number of single-character insertions and deletions to turn $X$
into $Y$. Then
\[
d(X,Y) = m + n - 2 \cdot \mathrm{LCS}(X,Y),
\]
where $\mathrm{LCS}(X,Y)$ is the length of the longest common subsequence.
\end{theorem}
\begin{proof}[Proof sketch]
Any common subsequence $S$ of length $\ell$ gives an edit script: delete the $m - \ell$
characters of $X$ not in $S$, then insert the $n - \ell$ characters of $Y$ not in $S$, for
$(m-\ell) + (n-\ell) = m + n - 2\ell$ operations. Maximising $\ell$ minimises the count, so
$d(X,Y) \le m + n - 2\,\mathrm{LCS}$. Conversely, any insert/delete script that converts $X$
to $Y$ leaves a set of untouched characters forming a common subsequence of length
$\ell'$, and uses at least $m+n-2\ell'$ operations, so $d(X,Y) \ge m+n-2\,\mathrm{LCS}$.
\end{proof}
Problem 5 asks you to verify this identity empirically by comparing an insert/delete-only
edit distance against $m + n - 2\,\mathrm{LCS}$ on random strings.

\section{Shortest Paths with Negative Edges: Bellman--Ford}

\subsection{Why Dijkstra is not enough}

Dijkstra's algorithm (Week 5) assumes all edge weights are non-negative; it commits to a
vertex's distance as soon as that vertex is extracted from the priority queue, which is only
valid when no later (negative) edge could improve it. With negative edges this assumption
fails, and Dijkstra can return wrong answers. Negative edges arise naturally -- e.g.\ in
\textbf{arbitrage detection} (currency exchange rates become edge weights after taking
logarithms, and a profitable cycle is a negative cycle) and in problems reduced to shortest
paths.

A subtlety: if a graph contains a \textbf{negative cycle} reachable on the way from $s$ to
$t$, the notion of ``shortest path'' is ill-defined -- one can loop around the cycle
arbitrarily many times to drive the cost to $-\infty$. So the right problem is: \emph{find
shortest paths assuming no negative cycle, and otherwise report that one exists.}

\subsection{The DP formulation}

The key idea is to bound the \emph{number of edges} a shortest path may use. In a graph
with no negative cycle, any shortest $s$--$v$ path is \emph{simple} (visits no vertex
twice), hence uses at most $n-1$ edges. So define
\[
\mathrm{OPT}(k, v) = \text{the minimum cost of an $s$--$v$ path using at most $k$ edges.}
\]
Reasoning about the \emph{last edge} of such a path:
\[
\mathrm{OPT}(k, v) = \min\Big( \mathrm{OPT}(k-1, v),\ \min_{(u,v)\in E}\big(\mathrm{OPT}(k-1,u) + w(u,v)\big) \Big),
\]
with $\mathrm{OPT}(0, s) = 0$ and $\mathrm{OPT}(0, v) = \infty$ for $v \ne s$. Either the
shortest path with $\le k$ edges actually uses $\le k-1$ edges (first term), or its last
edge is some $(u,v)$ and the prefix is an optimal $\le k-1$ edge path to $u$ (second term).
The true shortest-path distances are $\mathrm{OPT}(n-1, v)$.

\subsection{The algorithm}

In implementation we do not keep a separate row per $k$; we keep a single distance array
$\mathrm{dist}[\cdot]$ and \textbf{relax} every edge, repeating the whole pass $n-1$ times.
``Relaxing'' edge $(u,v)$ means: if $\mathrm{dist}[u] + w(u,v) < \mathrm{dist}[v]$, update
$\mathrm{dist}[v]$ (and record $u$ as $v$'s parent).

\begin{algorithm}[h]
\caption{Bellman--Ford (single-source shortest paths)}
\begin{algorithmic}[1]
\State $\mathrm{dist}[s] \gets 0$; \quad $\mathrm{dist}[v] \gets \infty$ for all $v \ne s$
\For{$k = 1$ to $n-1$}
    \For{each edge $(u, v, w) \in E$}
        \If{$\mathrm{dist}[u] + w < \mathrm{dist}[v]$}
            \State $\mathrm{dist}[v] \gets \mathrm{dist}[u] + w$; \quad $\mathrm{parent}[v] \gets u$
        \EndIf
    \EndFor
\EndFor
\State \Comment{Negative-cycle check: one more pass}
\For{each edge $(u, v, w) \in E$}
    \If{$\mathrm{dist}[u] + w < \mathrm{dist}[v]$}
        \State \Return ``negative cycle detected''
    \EndIf
\EndFor
\State \Return $\mathrm{dist}$ and $\mathrm{parent}$
\end{algorithmic}
\end{algorithm}

\begin{theorem}
After the $k$-th pass of edge relaxations, $\mathrm{dist}[v] \le \mathrm{OPT}(k,v)$ for
every $v$; in particular, if the graph has no negative cycle, after $n-1$ passes
$\mathrm{dist}[v]$ equals the true shortest-path distance from $s$ to $v$. The algorithm
runs in $O(mn)$ time.
\end{theorem}
\begin{proof}
By induction on $k$. Initially $\mathrm{dist}[s]=0=\mathrm{OPT}(0,s)$ and all others are
$\infty$. Suppose the claim holds after pass $k-1$. Consider any $v$ and an optimal
$\le k$-edge path $P$ to $v$. If $P$ uses $\le k-1$ edges, then by induction
$\mathrm{dist}[v]$ was already $\le \mathrm{OPT}(k-1,v) = \mathrm{cost}(P)$. Otherwise $P$
has a last edge $(u,v)$ and a prefix that is a $\le k-1$ edge path to $u$ of cost
$\mathrm{OPT}(k-1,u)$; by induction $\mathrm{dist}[u] \le \mathrm{OPT}(k-1,u)$ before pass
$k$, and when pass $k$ relaxes $(u,v)$ it sets $\mathrm{dist}[v] \le \mathrm{dist}[u] + w
\le \mathrm{cost}(P)$. Since shortest paths in a negative-cycle-free graph are simple and
use at most $n-1$ edges, $\mathrm{dist}[v]$ reaches the true distance after $n-1$ passes.
Each pass relaxes all $m$ edges, giving $O(mn)$ total.
\end{proof}

\subsection{Negative-cycle detection}

The correctness theorem says distances stabilise after $n-1$ passes \emph{when no negative
cycle exists}. This gives a clean test (Problem 7): run one \textbf{extra} ($n$-th) pass; if
\emph{any} edge can still be relaxed (some $\mathrm{dist}[u] + w < \mathrm{dist}[v]$), then a
shortest path of more than $n-1$ edges would be improving on a simple path, which is only
possible if a negative cycle is reachable. To recover the cycle itself, record the parent
pointer that caused the successful relaxation in the extra pass, then follow parent pointers
$n$ times (to guarantee landing inside the cycle) and read off the cycle. Problem 6
implements the distance computation and Problem 7 the detection.

\section{All-Pairs Shortest Paths: Floyd--Warshall}

\subsection{Problem and recurrence}

Sometimes we want shortest paths between \emph{every} pair of vertices. Running Bellman--Ford
from each source costs $O(mn^2)$; for dense graphs the \textbf{Floyd--Warshall} algorithm is
simpler and competitive at $O(n^3)$. Its DP is indexed by the set of \emph{intermediate
vertices} allowed on a path.

Number the vertices $1, \ldots, n$. Let $d_k(i,j)$ be the length of the shortest path from
$i$ to $j$ whose intermediate vertices all lie in $\{1, \ldots, k\}$. Considering whether an
optimal such path uses vertex $k$ as an intermediate:
\[
d_k(i,j) = \min\big(d_{k-1}(i,j),\ d_{k-1}(i,k) + d_{k-1}(k,j)\big),
\]
with base case $d_0(i,j) = w(i,j)$ if $(i,j) \in E$, $d_0(i,i) = 0$, and $\infty$ otherwise.
Either the optimal path avoids $k$ (first term), or it goes $i \rightsquigarrow k
\rightsquigarrow j$ through $k$ exactly once, splitting into two paths each using only
$\{1,\ldots,k-1\}$ as intermediates (second term). The final answer is $d_n(i,j)$.

\begin{algorithm}[h]
\caption{Floyd--Warshall (all-pairs shortest paths)}
\begin{algorithmic}[1]
\State $d[i][j] \gets w(i,j)$ for edges; $d[i][i] \gets 0$; $\infty$ otherwise
\For{$k = 1$ to $n$}
    \For{$i = 1$ to $n$}
        \For{$j = 1$ to $n$}
            \If{$d[i][k] + d[k][j] < d[i][j]$}
                \State $d[i][j] \gets d[i][k] + d[k][j]$
            \EndIf
        \EndFor
    \EndFor
\EndFor
\State \Return $d$
\end{algorithmic}
\end{algorithm}

As with Bellman--Ford, a single matrix is updated in place; the correctness of doing so
rests on the fact that the $k$-th outer iteration only \emph{reads} entries that are still
correct for the ``$\le k$ intermediates'' subproblem. A \textbf{negative cycle} reveals
itself as a strictly negative diagonal entry $d[i][i] < 0$ after the algorithm finishes.
Problem 8 implements Floyd--Warshall, and Problem 9 cross-checks it against Bellman--Ford:
on any graph with no negative cycle, Bellman--Ford from source $s$ must agree, vertex by
vertex, with row $s$ of the Floyd--Warshall matrix.

\section{Shortest and Longest Paths in a DAG}

\subsection{Topological order makes DP trivial}

In a \textbf{directed acyclic graph} there are no cycles at all -- in particular no negative
cycles -- so even longest paths are well-defined (unlike in general graphs, where
longest-simple-path is NP-hard). A topological order $v_1, v_2, \ldots, v_n$ (every edge
goes from earlier to later in the order; recall Week 3) lets us process vertices in an order
where, by the time we reach $v$, all of $v$'s predecessors are already finalised. The DP is
then a single sweep:
\[
\mathrm{dist}(v) = \min_{(u,v) \in E} \big( \mathrm{dist}(u) + w(u,v) \big),
\qquad \mathrm{dist}(s) = 0,
\]
processed in topological order. For \emph{longest} paths, replace $\min$ with $\max$ (and
initialise unreachable vertices to $-\infty$). Because each edge is examined once, the whole
computation is $\mathbf{O(n + m)}$ -- faster than both Dijkstra ($O(m \log n)$) and
Bellman--Ford ($O(mn)$), and it handles negative weights for free.

\begin{algorithm}[h]
\caption{DAG Shortest Path (relax in topological order)}
\begin{algorithmic}[1]
\State Compute a topological order of the vertices
\State $\mathrm{dist}[s] \gets 0$; \quad $\mathrm{dist}[v] \gets \infty$ for all $v \ne s$
\For{each vertex $u$ in topological order}
    \For{each edge $(u, v, w)$ leaving $u$}
        \If{$\mathrm{dist}[u] + w < \mathrm{dist}[v]$}
            \State $\mathrm{dist}[v] \gets \mathrm{dist}[u] + w$; \quad $\mathrm{parent}[v] \gets u$
        \EndIf
    \EndFor
\EndFor
\State \Return $\mathrm{dist}$ and $\mathrm{parent}$
\end{algorithmic}
\end{algorithm}

This is the right tool when a problem has an inherent acyclic structure -- task scheduling
along dependencies, the critical path in a project (a \emph{longest} path), or counting
paths in a grid. Problem 10 implements DAG shortest paths and checks them against the more
general Bellman--Ford, confirming that on a DAG the simpler linear-time sweep suffices where
Dijkstra-style machinery would be overkill.

\section{Two More Classic 2-D Dynamic Programs}

\subsection{Matrix chain multiplication}

Given a chain of matrices $A_1 A_2 \cdots A_n$ with $A_i$ of dimension $p_{i-1} \times
p_i$, matrix multiplication is associative, so any parenthesisation yields the same product
-- but the \emph{cost} (number of scalar multiplications) varies dramatically. The DP is
indexed by sub-chains: let $M(i,j)$ be the minimum cost of multiplying $A_i \cdots A_j$.
Reasoning about the \emph{outermost} multiplication (the last split point $k$):
\[
M(i,j) = \min_{i \le k < j} \big( M(i,k) + M(k+1, j) + p_{i-1}\,p_k\,p_j \big),
\qquad M(i,i) = 0.
\]
There are $O(n^2)$ subproblems, each minimising over $O(n)$ split points, for $O(n^3)$ time.
Recording the optimal split $k$ for each $(i,j)$ lets us reconstruct the optimal
parenthesisation as a string (Problem 11). This is the textbook example of DP over
\emph{contiguous sub-intervals}, a third indexing pattern beyond prefixes and vertex sets.

\subsection{Grid paths}

Consider an $r \times c$ grid where you may move only \emph{right} or \emph{down}. Two
classic questions: \emph{(a)} how many distinct paths lead from the top-left to the
bottom-right corner, and \emph{(b)} if each cell has a cost, what is the
\emph{minimum-cost} path? Both are one-line recurrences. For counting,
\[
P(i,j) = P(i-1,j) + P(i,j-1), \qquad P(0,0) = 1,
\]
and for minimum path sum over cell costs $c(i,j)$,
\[
S(i,j) = c(i,j) + \min\big(S(i-1,j),\, S(i,j-1)\big),
\]
each filled in $O(rc)$ time. The grid is implicitly a DAG (every move increases $i+j$), so
these are just the DAG recurrences specialised to a grid; an optimal path is recovered by
the usual traceback (Problem 12).

\section{Summary and Looking Ahead}

This week extended dynamic programming from one-dimensional to two-dimensional (and
interval- and vertex-set-indexed) subproblems, through several fundamental algorithms:

\begin{itemize}
    \item \textbf{Sequence alignment / edit distance}: the $O(mn)$ two-dimensional
    recurrence over prefix pairs, with $O(m+n)$ traceback for the alignment, and
    Hirschberg's $O(m+n)$-space refinement.
    \item \textbf{Longest common subsequence}: a sibling recurrence, related to
    insert/delete edit distance by $d = m + n - 2\,\mathrm{LCS}$.
    \item \textbf{Bellman--Ford}: shortest paths with negative edges in $O(mn)$ via a DP on
    the number of edges used, plus a one-extra-pass test for negative cycles.
    \item \textbf{Floyd--Warshall}: all-pairs shortest paths in $O(n^3)$ via a DP on the set
    of allowed intermediate vertices.
    \item \textbf{DAG shortest/longest paths}: a single $O(n+m)$ sweep in topological order,
    handling negative weights and longest paths for free.
    \item \textbf{Matrix chain multiplication} and \textbf{grid paths}: DP over sub-intervals
    and over a grid, both with reconstruction.
\end{itemize}

The unifying lesson is the DP recipe: define subproblems, reason about the \emph{last
decision} to get a recurrence, order the computation so dependencies come first, and trace
back to reconstruct an actual optimal solution.

Next week we turn to \textbf{network flow} (Kleinberg \& Tardos, Ch.\ 7): the maximum-flow /
minimum-cut problem and the Ford--Fulkerson method. Flow is a different algorithmic paradigm
from DP, but it shares this week's graph-theoretic setting and -- as we will see -- a great
many seemingly unrelated problems (bipartite matching, image segmentation, project
selection) reduce to it.

\end{document}
